\chapter{SWEB: putting it together}
\label{ch:sweb}

SWEB is the course's teaching kernel: small enough to read, real enough to
boot on x86-64 hardware (in QEMU). This last chapter is a map --- where
things are, how the machine gets from power-on to the shell, where each
idea of chapters 1--10 lives in the code --- and a guide to the A1 work:
what changes, and how to test it so that it survives the tenth run.
\module{11}'s worksheet walks through all of it with questions.

\section{The source tree}

\begin{center}
\begin{tabular}{L{5.4cm}L{8.4cm}}
\toprule
\textbf{Path} & \textbf{Contents} \\
\midrule
\file{common/source/kernel/} & \code{main.cpp} (\code{startup}), \code{Thread},
\code{UserProcess}, \code{Scheduler}, \code{Syscall}, \code{Loader},
\code{ProcessRegistry}, locks (\code{Lock}, \code{Mutex}, \code{SpinLock},
\code{Condition}), \code{CleanupThread}, \code{IdleThread} \\
\file{common/source/mm/} & \code{PageManager} (frames),
\code{KernelMemoryManager} (kernel heap), \code{PageFaultHandler} \\
\file{common/source/fs/}, \file{console/} & VFS, minix file system, terminal,
\code{kprintf} and \code{debug()} \\
\file{common/include/kernel/} & headers, \code{syscall-definitions.h} (numbers,
shared with the libc), \code{user\_progs.h} \\
\file{common/include/console/debug.h} & debug categories (switch
output on/off per category) \\
\file{arch/x86/64/} & page tables (\code{ArchMemory}), registers and context
switch (\code{ArchThreads}, \code{arch\_interrupts.S}), the IDT and all
interrupt handlers (\code{InterruptUtils.cpp}), \code{offsets.h}
(\code{USER\_BREAK}) \\
\file{userspace/libc/} & SWEB's own small C library: \code{write.c},
\code{pthread.c} (stubs), \code{nonstd.c} (\code{\_start}) \\
\file{userspace/tests/} & user programs; every \code{x.c} becomes
\code{x.sweb} on the disk image \\
\bottomrule
\end{tabular}
\end{center}

\section{From power-on to the shell}

\begin{enumerate}
  \item GRUB loads the kernel (the menu has no timeout --- press Enter);
    early assembly sets up paging and jumps to \code{startup()}.
  \item \code{startup()} creates the console, the scheduler (with the
    cleanup and idle threads), mounts the root file system and adds the
    \code{ProcessRegistry} thread; then it enables interrupts and yields ---
    from now on the scheduler runs everything.
  \item \code{ProcessRegistry::Run} mounts the user-program partition at
    \file{/usr} and creates a \code{UserProcess} for each entry of
    \code{user\_progs} --- by default only \file{/usr/shell.sweb}.
  \item The \code{UserProcess} constructor opens the binary, creates a
    \code{Loader} (reads the ELF headers), maps one stack page below
    \code{USER\_BREAK}, sets the user registers (\code{rip} = ELF entry =
    \code{\_start}, \code{rsp} = \code{USER\_BREAK - 8}) and the address
    space; the scheduler switches to it in user mode.
  \item Every first touch of code or data page-faults and is loaded by
    \code{Loader::loadPage}. The shell reads commands with \code{sc\_read}
    and starts programs with \code{createprocess}.
\end{enumerate}

\section{Every chapter in the code}

\begin{center}
\begin{tabular}{L{2.2cm}L{11.6cm}}
\toprule
\textbf{Chapter} & \textbf{Where in SWEB} \\
\midrule
1 processes & \code{UserProcess} (= a thread today), \code{ProcessRegistry},
\code{createprocess}; no \code{fork}/\code{exec} yet (Team F) \\
2 threads & \code{pthread.c} stubs; \code{Thread} with \code{kernel\_stack\_}
and two register sets \\
3 races & a preemptive kernel: interrupts are on during system calls \\
4 locks & \code{SpinLock} (test-and-set + yield), \code{Mutex} (sleeps),
\code{ArchThreads::testSetLock}, interrupts off in \code{schedule()} \\
5 condvars & \code{Condition}; \code{Scheduler::sleep}/\code{wake} with the
wait-until-sleeping loop \\
6 deadlocks & \code{holding\_lock\_list\_}, \code{lock\_waiting\_on\_},
\code{doChecksBeforeWaiting}; no locks in interrupt handlers \\
7 syscalls & \code{int \$0x80}, \code{arch\_syscallHandler},
\code{Syscall::syscallException}, \code{buffer + size > USER\_BREAK} checks \\
8 memory & \code{ArchMemory::mapPage/unmapPage}, \code{PageManager},
\code{PageFaultHandler}, \code{Loader::loadPage}, identity mapping \\
9 scheduling & round robin in \code{Scheduler::schedule}, one tick per quantum,
\code{yield} = \code{int \$65} \\
10 threads & \code{Thread} states, \code{kill} $\to$ \code{ToBeDestroyed} $\to$
cleanup thread \\
\bottomrule
\end{tabular}
\end{center}

\section{What A1 changes}

\paragraph{Phase 0 (all four, 1--4 Oct): split process from thread.} A
process owns the address space (\code{Loader}/\code{ArchMemory}), the file
descriptors and the PID; a thread owns its stacks, registers, state and
TID. Consequences to design together: how a thread finds its process
(a pointer; keeping \code{Thread::loader\_} as a non-owning shortcut
saves many edits), who destroys the process and when (after its
\emph{last} thread, in the cleanup thread --- never on a stack or in an
address space that is about to disappear), how ids are handed out (today
every \code{tid\_} is 0), and what \code{ProcessRegistry} counts.

\paragraph{Team P.} \code{pthread\_create} (P1): the system call, argument
checks, a new thread with its own user stack and a start wrapper.
\code{pthread\_exit} (P2): the return value survives the thread, the
stack is freed, the process keeps running. \code{pthread\_join} (P1): block
until the target has finished, receive exactly the stored value; a
detached thread cannot be joined. \code{pthread\_cancel} (P2): end the
target at a defined cancellation point.

\paragraph{Team F.} \code{fork} with copy-on-write, \code{exec} --- and,
together with Team P, their behaviour in multi-threaded processes
(\code{fork} duplicates only the calling thread; \code{exit} tears down
every thread).

\section{Testing}

The course grades with your own tests and with the Tacos test system;
neither replaces the other. Practical rules:

\begin{itemize}
  \item Write the test program first (\file{userspace/tests/}); re-run
    CMake after adding one, and build twice if the new program is missing
    from the image.
  \item Test the edges, not only the happy path: bad user pointers,
    joining yourself or twice, main exiting first, a thread creating
    threads, many threads, create/join in a loop (leaks).
  \item Run every test many times: \module{11}'s \code{sweb\_run.py -n 10}
    boots SWEB headless ten times and reports hangs (exit 1) and kernel
    panics (exit 2). Races appear on the tenth run.
  \item Use the debug categories (\code{debug.h}) instead of scattering
    \code{kprintf}; switch noisy ones off. For hard cases: \code{make
    qemugdb} + \code{make rungdb}.
  \item Assertions are cheap documentation: \code{assert(lock.isHeldBy(currentThread))}
    at the top of a function that needs the lock.
\end{itemize}

\section{Pitfalls that cost days}

\begin{itemize}
  \item \code{new}/\code{delete} with interrupts off or in a handler: the
    kernel-heap lock may be held by the interrupted thread.
  \item Sleeping while holding a lock the waker needs; \code{Condition::wait}
    asserts against extra held locks for a reason.
  \item Deleting a thread (or an address space) that is still running or
    still in use --- defer to the cleanup thread.
  \item Trusting user pointers, or checking them with an overflowing sum.
  \item Busy-waiting (\code{while (!done) yield();}) instead of sleeping on a
    condition.
  \item Two threads faulting on the same missing page table at once:
    \code{mapPage} needs a lock once a process has several threads.
  \item Checking once, using later (TOCTOU) --- in user memory and in
    kernel state: ``is the thread still alive?'' must be answered and acted
    on under the same lock.
\end{itemize}

\begin{keybox}
\begin{itemize}
  \item Know the map: \file{common/source/kernel}, \file{mm},
    \file{arch/x86/64}, \file{userspace/libc}, \file{userspace/tests}.
  \item Boot path: \code{startup} $\to$ \code{ProcessRegistry} $\to$
    \code{UserProcess} $\to$ lazy loading by page faults.
  \item A1 starts with the process/thread split --- agree on ownership,
    lifetime, ids and locking as a group.
  \item Test the edges, many times, headless; read the debug output.
\end{itemize}
\end{keybox}
